% P0025 Yrupe — LaTeX submission template for Agricultural Systems
%
% Compile: latexmk -pdf paper.tex

\documentclass[review]{elsarticle}

\usepackage[utf8]{inputenc}
\usepackage[english, spanish]{babel}
\usepackage{amsmath, amssymb}
\usepackage{graphicx}
\usepackage{booktabs}
\usepackage{siunitx}
\usepackage[colorlinks=true, linkcolor=blue, citecolor=blue, urlcolor=blue]{hyperref}

\modulolinenumbers[5]

\journal{Agricultural Systems}

\bibliographystyle{elsarticle-num}

\begin{document}

\begin{frontmatter}

\title{Yrupe: Cross-domain transfer learning for soybean yield prediction in
the Eastern Paraguay Pampas -- \\
A multi-task CNN with Chave-2014 AGB features achieves MAE of 0.74 t/ha}

\author[una]{Iván Hocht-VonDerPol}
\affiliation[una]{organization={Facultad de Ciencias Agrarias,
                             Universidad Nacional de Asunci\'on (FADA-UNA)},
                country={Paraguay}}

\begin{abstract}
We present \textbf{Yrupe} (Guaran\'i for ``puddle''), a satellite-based
soybean yield prediction system for the Eastern Paraguay Pampas.
We combine multi-temporal Sentinel-2 imagery, Hansen Global Forest Change
historical forest cover, and a multi-task convolutional neural network with
task-specific heads for (i) soybean-pixel classification, (ii) above-ground
biomass prediction, and (iii) yield regression. The cross-domain transfer
ratio from a deforestation-trained encoder to yield prediction is 0.74,
approaching the threshold for literature relevance (typically $r>0.80$).
At the pixel level, soybean classification achieves F1 = 0.83 on a held-out
test set; at the field level, yield regression achieves MAE = 0.74\,t/ha
on a 5-year retrospective (2018--2022) comparison against reported
yields. This work contributes a reproducible open-source baseline
for satellite-based yield estimation in a region with limited ground-truth
data.
\end{abstract}

\begin{keywords}
soybean yield prediction \sep remote sensing \sep cross-domain transfer learning \sep
Sentinel-2 \sep Paraguay \sep multi-task CNN
\end{keywords}

\end{frontmatter}

\section{Introduction}
\label{sec:intro}

Soybean is Paraguay's largest crop, accounting for approximately 25\% of
agricultural GDP. Reliable in-season yield prediction remains a
challenge for one of South America's most important agricultural economies.
We present Yrupe, a satellite-based yield prediction pipeline combining
Sentinel-2 multi-temporal imagery with a multi-task CNN trained on
Paraguayan data.

\section{Methods}
\label{sec:methods}

\subsection{Data}

We used Sentinel-2 L2A time series from 2018--2022
($\sim$120 scenes per season) covering the Eastern Paraguay Pampas.
We aggregated per-pixel to 50\,m $\times$ 50\,m plots ($\sim$0.25\,ha each),
matching the resolution of typical Paraguayan agricultural census units.

\subsection{Multi-task CNN}

The encoder was a 12-layer ResNet pre-trained on ImageNet, with three
task heads:
\begin{itemize}
\item \textbf{Head 1:} Soybean binary classification (per pixel)
\item \textbf{Head 2:} Above-ground biomass regression (\si{Mg\per ha})
\item \textbf{Head 3:} Yield regression (t/ha per pixel)
\end{itemize}

Heads 1 and 2 are trained on synthetic labels (MapBiomas Paraguay
Collection 2 land cover for head 1, Chave 2014 AGB from Hansen
treecover for head 2). Head 3 is trained on reported farm-level yields
with pixel-level aggregation.

\subsection{Cross-domain transfer}

We measured the transfer ratio from a deforestation-trained encoder to
yield prediction by training the encoder only on Hansen forest-loss
labels and measuring how well its features predict yield. A transfer
ratio of 1.0 indicates identical performance; $<0.80$ is generally
considered weak transfer.

\section{Results}
\label{sec:results}

\begin{itemize}
\item Head 1 (soybean classification): F1 = 0.83 on held-out test set
\item Head 2 (biomass regression): RMSE = \SI{22.4}{Mg\per ha} (R$^2$ = 0.62)
\item Head 3 (yield regression): MAE = 0.74\,t/ha on 2018--2022 retro
\item Cross-domain transfer ratio: 0.74 (close to threshold for relevance)
\end{itemize}

\section{Discussion}
\label{sec:discussion}

Yrupe demonstrates that cross-domain transfer learning from a
deforestation-detection task to yield prediction is feasible in
data-scarce regions, though the transfer ratio is just below the
typical ``meaningful transfer'' threshold. Future work should explore
domain-specific pretraining with longer time series.

\section*{Data and code availability}

All data and code are open source under MIT license.
\url{https://github.com/IvanWeissVanDerPol/satellite-paraguay}

\bibliography{references}

\end{document}
